\section{Synth Modules}\label{synth-modules}

A closer look at the synthesis modules used across the different voices.

\subsection{Oscillator}\label{oscillator}

The oscillators provide 6 different waveforms:

\begin{itemize}
\tightlist
\item
  Sine
\item
  Triangle
\item
  Saw
\item
  Rectangle
\item
  Noise
\item
  Crash
\end{itemize}

The first 5 are classical analogue waveforms realised using bandlimited
wavetable oscillators. The 6th is a TR-909-style crash sample ---
included as a bonus, since a realistic crash is nearly impossible to
synthesise using the available voice structure.

Any samples loaded will be listed after the stock waveforms: one-shot
samples first, then loops, alphabetically by the name of the file they
were imported from.

\subsubsection{Wavetable Interpolation}\label{wavetable-interpolation}

If the oscillator waveform is modulated, interpolation between waveforms
can be switched on in the global settings menu, producing a smoother
transition between waves. The interpolation effect applies only to the
main oscillators of the voices, and is only performed on one oscillator
at a time to limit CPU overhead. However, the calculation is performed
dynamically as needed. That is, one wave-modulated oscillator will
always be interpolated, two wave-modulated oscillators will be
interpolated provided the modulation is \emph{slow enough}, etc.

This is particularly useful when importing a series of
progressively-changing single-cycle looped waveforms of the same length
in a sequence as part of the `LOOPS' folder ample import. Modulating the
waveform with interpolation on when the oscillator is set to play one of
the samples produces a smooth timbral change between adjacent samples.

\subsection{Filter}\label{filter}

The filter is a 2-pole (12 dB) state variable filter (SVF) realised as a
non-linear zero-delay-feedback model. It shapes the harmonic content of
the sound --- for example, a hihat needs a highpass filter to remove low
frequencies, while a bandpass is useful for a clap.

There are 3 parameters that control the filter response: frequency,
resonance, and drive.

\subsubsection{Frequency}\label{frequency}

The frequency determines the operating point of the filter in the
spectrum. A lowpass filter cuts all frequencies above the set cutoff
frequency.

\subsubsection{Resonance}\label{resonance}

The resonance controls the feedback path of the filter. Higher resonance
settings amplify the frequencies around the operating point (i.e.~the
cutoff frequency) more and more.

\subsubsection{Drive}\label{drive}

Controls how hard the filter is driven. More drive yields more
distortion. At drive 0 the filter is clean and nearly linear. With
higher levels, a soft clipper and the slew-rate limit of the integrators
engage. Low settings only affect resonance peaks; high settings distort
the whole signal. Because the soft clipper scales down excessive peaks,
audible resonance is reduced at high drive settings.

\subsubsection{Filter Types}\label{filter-types}

There are several filter types, each with its own characteristics. The
plots below show each filter with 3 resonance settings (low = red,
medium = green, high = blue). The cutoff frequency is identical in every
plot, marked by a vertical dashed line.

\textbf{Lowpass} --- Removes high frequencies. All frequencies above the
cutoff are reduced gradually.

\begin{figure}[H]
  \centering
  \includegraphics[width=0.85\linewidth]{filter_lowpass.png}
  \caption{Lowpass filter response.}
\end{figure}

\textbf{Highpass} --- Removes low frequencies. All frequencies below the
cutoff are reduced gradually.

\begin{figure}[H]
  \centering
  \includegraphics[width=0.85\linewidth]{filter_highpass.png}
  \caption{Highpass filter response.}
\end{figure}

\textbf{Bandpass} --- Removes frequencies both above and below the
cutoff.

\begin{figure}[H]
  \centering
  \includegraphics[width=0.85\linewidth]{filter_bandpass.png}
  \caption{Bandpass filter response.}
\end{figure}

\textbf{Unit Gain Bandpass} --- A scaled bandpass where gain is always
adjusted to a maximum of 0 dB. Unlike the normal bandpass, resonance
controls the width of the passband rather than the peak amplitude.

\begin{figure}[H]
  \centering
  \includegraphics[width=0.85\linewidth]{filter_unitgain_bandpass.png}
  \caption{Unit gain bandpass filter response.}
\end{figure}

\textbf{Notch} --- Removes frequencies around the cutoff. Resonance
controls the width of the stopband.

\begin{figure}[H]
  \centering
  \includegraphics[width=0.85\linewidth]{filter_notch.png}
  \caption{Notch filter response.}
\end{figure}

\textbf{Peak} --- Amplifies frequencies around the cutoff; other
frequencies pass nearly unaltered. Resonance controls the amount of
amplification.

\begin{figure}[H]
  \centering
  \includegraphics[width=0.85\linewidth]{filter_peak.png}
  \caption{Peak filter response.}
\end{figure}

\subsection{Envelopes}\label{envelopes}

Envelopes generate a varying control signal to modulate other synthesis
parameters. Whenever a voice is triggered, the envelope restarts. The
signal rises at the speed set by attack, reaches maximum amplitude, then
falls back to zero at the speed set by decay.

\begin{figure}[H]
  \centering
  \includegraphics[width=0.7\linewidth]{envelope_linear.png}
  \caption{Linear envelope curve (slope = 63).}
\end{figure}

\subsubsection{Attack / Decay Times}\label{attack-decay-times}

These 2 parameters control the duration of each stage. Higher values =
longer times. If attack is set to zero, the envelope starts directly in
the decay stage.

\subsubsection{Slope}\label{slope}

The slope parameter controls the curve shape. A value of 63 gives a
linear envelope. Lower values produce an exponential curve; higher
values produce a logarithmic curve.

\begin{figure}[H]
  \centering
  \includegraphics[width=0.7\linewidth]{envelope_exponential.jpeg}
  \caption{Low slope setting — exponential decay stage.}
\end{figure}

\begin{figure}[H]
  \centering
  \includegraphics[width=0.7\linewidth]{envelope_logarithmic.jpeg}
  \caption{High slope setting — logarithmic decay stage.}
\end{figure}

\subsubsection{Repeat}\label{repeat}

The repeat feature simulates clap and other sounds with a fuzzy or
rattling attack. When the repeat count is set above zero, the slow
attack stage is replaced by a repeat stage in which the envelope
retriggered rapidly. The attack parameter controls the duration of the
complete repeat phase.

\begin{figure}[H]
  \centering
  \includegraphics[width=0.85\linewidth]{envelope_repeat.png}
  \caption{Envelope in repeat mode (count = 3). Blue = repeat phase, red = decay phase.}
\end{figure}

\subsection{Transient Generator}\label{transient-generator}

The transient generator shapes the very beginning of a sound. Because
the first few milliseconds heavily influence how a sound is perceived,
this is a powerful tool for varying the attack character.

\subsubsection{\texorpdfstring{Snappy Mode
(\texttt{Snp})}{Snappy Mode (Snp)}}\label{snappy-mode-snp}

An additional fast pitch envelope modulates the oscillator pitch,
producing a snap or click. The \param{vol} parameter controls modulation
depth; \param{frq} controls the decay time.

\subsubsection{\texorpdfstring{Offset Mode
(\texttt{Off})}{Offset Mode (Off)}}\label{offset-mode-off}

Adjusts the oscillator start phase. At volume 0, oscillators reset to a
zero crossing on each trigger. At maximum volume, the waveform starts at
its highest amplitude, generating a loud pop.

\subsubsection{Sample Mode}\label{sample-mode}

A short (\textasciitilde50 ms) 8-bit ROM sample is mixed into the voice
output whenever it is triggered. Samples are hardcoded into the firmware
and cannot be changed by the user.

Available transient samples:

\begin{paramtable}
Snp & Snappy Mode     & (Pitch envelope — see above.) \\
Off & Offset Mode     & (Phase offset — see above.) \\
Clk & Click           & Short click transient. \\
Ck2 & Click 2         & Alternate click. \\
Tik & Tick            & A tick-style transient. \\
Kik & Acoustic kick   & An acoustic kick drum sample. \\
Rim & Rim shot        & A rim shot sample. \\
Drp & Drip            & A dripping sound. \\
Hat & Hat             & Sampled hi-hat. \\
Clp & Clap            & 808-style clap. \\
Ki2 & Kick 2          & Second kick sample. \\
Sna & Snare           & A snare sample. \\
Tom & Tom             & A tom sample. \\
Sp2 & Snap            & A finger snap. \\
\end{paramtable}

\subsection{Sample Rate Reduction}\label{sample-rate-reduction}

The sample rate reducer gradually reduces the sample rate from 44 kHz
downward. Reducing the rate gives a lo-fi character. The further it is
reduced, the more odd harmonics are added as overtones fold down at the
Nyquist frequency --- until the sound is completely destroyed.

\subsection{Distortion}\label{distortion}

The distortion unit on the mixer page is a variable waveshaper. It
ranges from no distortion through soft saturation up to hard clipping.

\begin{figure}[H]
  \centering
  \includegraphics[width=0.7\linewidth]{distortion_curve.jpeg}
  \caption{Distortion waveshaper transfer curves: low (blue), mid (green), max (red).}
\end{figure}

\subsection{LFO}\label{lfo}

Each voice has a low frequency oscillator. LFOs are identical to the
audio oscillators but run at a lower maximum rate. They are not audible
directly --- they modulate other parameter values over time.

Example: use a sine LFO to slowly open and close a voice's filter
cutoff.

\subsubsection{Setting Up Modulations}\label{setting-up-modulations}

The modulation target is selected on the second page of the LFO menu:

\oled{rtg \quad off \quad voi \quad Dst <}{off \hspace{1.5em} 0 \hspace{2.2em} 1 \hspace{2.5em} coa}

The target is specified by 2 parameters: \textbf{Voice} and
\textbf{Destination}. Changing the voice parameter updates the
destination list to show all available parameters of the selected voice.

\attention{The LFO modulation scheme differs from most synthesisers. The destination value is
\emph{multiplied} by the LFO value, meaning it is always modulated between its current value and
zero. It can never exceed the original parameter value. \textbf{Modulating a parameter set to
zero has no effect.}}

The LFO can be synced to the sequencer clock --- at \(\frac{1}{4}\) note
sync, the LFO cycles exactly 4 times per bar. When sync is active, the
frequency parameter is ignored.

\subsubsection{Using The LFO as an
Envelope}\label{using-the-lfo-as-an-envelope}

You can use the LFO as an additional envelope by enabling the retrigger
feature and choosing one of the special `one-shot' waveforms. Set the
retrigger source to the same voice as the modulation destination so the
LFO will restart on every played note (this does not need to be the same
voice as the one the LFO is on). Additional `one-shot' LFOs are listed
in the LFO waveforms as si1, tr1, etc. to use with this; all the
waveforms with a `1' suffix generate their modulation envelope once when
triggered and then stop until triggered again, instead of cycling as the
other waveforms do. When using one of these waveforms, `offset' provides
an adjustable delay before the start of the waveform instead of
adjusting the waveform phase.

\subsection{Velocity Modulation}\label{velocity-modulation}

Note-on velocity is an important modulation source for bringing life to
static patches. Velocity can modulate 2 destinations:

\begin{itemize}
\tightlist
\item
  The voice volume.
\item
  Any single voice parameter.
\end{itemize}

\tip{By default, velocity modulates the voice volume. This hard-wired connection can be turned
off with the \param{vol} on/off parameter in the modulation menu. Disabling volume modulation
frees velocity for use as a third automation track.}

\subsection{Output Routing}\label{output-routing}

Each voice can be freely routed to any of the 4 output jacks. The jacks
operate either as 4 individual mono outputs (pan parameter has no
effect) or as 2 stereo pairs.

\tip{The synth detects whether a cable is inserted in each jack. If a voice is routed to an
unconnected output, it automatically plays on the nearest connected output instead.}
